
This work investigates whether residual-corrected Jacobian stable rank can be controlled during transformer pretraining via gradient-based regularization. The approach combined Hutchinson trace estimation (10 Rademacher probes) for Frobenius norm and randomized power iteration (5 iterations) for spectral norm, differentiated through PyTorch autodiff. Experiments on GPT-2 (125M parameters) trained on C4 for 5000 steps produced zero stable rank measurements across all 12 layers, perplexity of 45,792 (baseline: 59.3), and regularization losses of -17.5 billion. Root cause analysis identified three failure modes: insufficient Hutchinson probe count for 768-dimensional embeddings, power iteration non-convergence in residual-corrected computation graphs, and gradient pathologies in deep autodiff chains. The baseline model converged successfully (perplexity 59.3), isolating the failure to the regularization mechanism. This negative result documents specific numerical instabilities in gradient-based spectral control and suggests post-hoc SVD-based observational studies as an alternative approach.
